% © 2025 Ing. David Jaroš — CC BY-NC-ND 4.0
%
% This work is licensed under a Creative Commons Attribution-NonCommercial-NoDerivatives
% 4.0 International License (CC BY-NC-ND 4.0).

\documentclass[11pt,a4paper]{article}
\usepackage{amsmath, amssymb, amsthm}
\usepackage{hyperref}
\usepackage{geometry}
\geometry{margin=2.5cm}

\newtheorem{axiom}{Axiom}
\newtheorem{assumption}{Assumption}
\newtheorem{remark}{Remark}

\title{Canonical Axioms\\
\large Unified Biquaternion Theory — Canonical Theory Series}
\author{Ing.\ David Jaroš}
\date{2025}

\begin{document}
\maketitle

\begin{abstract}
This document lists the minimal set of axioms (foundational assumptions) of
the Unified Biquaternion Theory (UBT). These axioms are not derived — they
are the starting point from which all UBT results are derived. Each axiom is
stated precisely and cross-referenced to the primary source file where it is
used and discussed. This list is stable (Layer 0) and should not be modified
without explicit author instruction.
\end{abstract}

\section*{Notation}

\begin{tabular}{ll}
$\mathbb{H}$ & real quaternion algebra \\
$\mathbb{C}$ & complex numbers \\
$\mathcal{B} = \mathbb{H}\otimes_\mathbb{R}\mathbb{C}$ & biquaternion algebra \\
$\tau = t + i\psi$ & complex time ($t$ real, $\psi$ imaginary component) \\
$\Theta(q,\tau)$ & fundamental biquaternionic field \\
$\kappa$ & coupling constant (gravitational analogue) \\
$\mathcal{T}$ & source tensor (energy-momentum) \\
$\nabla^\dagger$ & biquaternionic conjugate gradient operator \\
\end{tabular}

\bigskip

% ------------------------------------------------------------------
\section{Layer 0: Structural Axioms (Exact, No Free Parameters)}
% ------------------------------------------------------------------

These axioms define the algebraic and geometric structure.
Their status is \textbf{[L0]} — exact by construction.

\begin{axiom}[Biquaternion Algebra]\label{ax:algebra}
The fundamental mathematical arena is the biquaternion algebra
$\mathcal{B} = \mathbb{H}\otimes_\mathbb{R}\mathbb{C}$, which is isomorphic
to the algebra of $2\times 2$ complex matrices:
\begin{equation}
  \mathcal{B} \;\cong\; \mathrm{Mat}(2,\mathbb{C}).
\end{equation}
\emph{Source}: \texttt{canonical/fields/biquaternion\_algebra.tex}, \S1;
\texttt{canonical/algebra/involutions\_Z2xZ2xZ2.tex}.
\end{axiom}

\begin{axiom}[Fundamental Field]\label{ax:field}
The fundamental dynamical object is a biquaternionic field:
\begin{equation}
  \Theta \;:\; \mathbb{R}^{1,3}\times\mathbb{C} \;\longrightarrow\; \mathcal{B},
  \qquad
  (q,\tau) \;\mapsto\; \Theta(q,\tau),
\end{equation}
where $q\in\mathbb{R}^{1,3}$ is the real spacetime coordinate and
$\tau = t + i\psi$ is complex time.
All physical fields (metric, gauge fields, matter) are encoded in $\Theta$.
\emph{Source}: \texttt{canonical/fields/theta\_field.tex}, \S1.
\end{axiom}

\begin{axiom}[Complex Time]\label{ax:complextime}
Time is extended to the complex plane:
\begin{equation}
  \tau \;=\; t + i\psi, \qquad t,\psi \in \mathbb{R}.
\end{equation}
The imaginary component $\psi$ parametrizes an internal phase (fiber) direction.
Real observables correspond to the $\psi=0$ section.
\emph{Source}: \texttt{canonical/fields/biquaternion\_time.tex}, \S1;
\texttt{THEORY/math/fields/biquaternion\_time.tex}.
\end{axiom}

% ------------------------------------------------------------------
\section{Layer 1: Dynamical Axioms (Proved Given Layer 0)}
% ------------------------------------------------------------------

\begin{axiom}[Field Equation]\label{ax:fieldeq}
The field $\Theta$ satisfies the biquaternionic wave equation:
\begin{equation}\label{eq:tshirt}
  \nabla^\dagger \nabla\,\Theta(q,\tau) \;=\; \kappa\,\mathcal{T}(q,\tau),
\end{equation}
where $\nabla^\dagger\nabla$ is the biquaternionic Laplacian and $\mathcal{T}$
is the biquaternionic source tensor.
This equation is called the ``T-shirt formula'' and is the master equation of UBT.
\emph{Source}: \texttt{canonical/fields/theta\_field.tex}, \S3.
\end{axiom}

\begin{axiom}[Real Projection]\label{ax:projection}
Physical observables are recovered by projecting the biquaternionic field
to real-valued quantities:
\begin{equation}
  g_{\mu\nu}(x) \;=\; \mathrm{Re}\bigl(\mathcal{G}_{\mu\nu}[\Theta]\bigr),
\end{equation}
where $\mathcal{G}_{\mu\nu}$ is the biquaternionic metric tensor derived from $\Theta$.
\emph{Source}: \texttt{canonical/geometry/metric.tex}, \S2;
\texttt{canonical/geometry/phase\_projection.tex}.
\end{axiom}

\begin{axiom}[Action Principle]\label{ax:action}
The dynamics of $\Theta$ are governed by the variational principle
$\delta S = 0$ with the UBT action:
\begin{equation}
  S[\Theta] \;=\; \int \mathrm{d}^4x\,\sqrt{-g}\,
  \Bigl[\,
    \frac{1}{2\kappa}\,R[\Theta]
    \;+\;
    \mathcal{L}_{\mathrm{matter}}[\Theta]
    \;+\;
    \mathcal{L}_{\mathrm{gauge}}[\Theta]
  \,\Bigr],
\end{equation}
where $R[\Theta]$ is the Ricci scalar derived from the projected metric,
$\mathcal{L}_{\mathrm{matter}}$ is the matter Lagrangian, and
$\mathcal{L}_{\mathrm{gauge}}$ is the gauge field Lagrangian.
\emph{Source}: \texttt{THEORY/canonical/canonical\_action.tex};
\texttt{unified\_biquaternion\_theory/ubt\_appendix\_1\_biquaternion\_gravity.tex}, \S1.
\end{axiom}

% ------------------------------------------------------------------
\section{Notes and Disclaimers}
% ------------------------------------------------------------------

\begin{remark}[On Consciousness and Speculative Extensions]
Axioms A1--A6 define the physical theory. Extensions to consciousness
(psychons, Complex Consciousness Theory) are speculative and are
maintained in \texttt{speculative\_extensions/} separately from the
core theory defined here.
\end{remark}

\begin{remark}[On Layer Numbering]
Layer 0 results are algebraic identities holding in $\mathcal{B}$ regardless
of physical interpretation. Layer 1 results are proved given the axioms above.
Layer 2 results require additional lemmas not yet proved.
This classification is used throughout the canonical bridges and audit documents.
\end{remark}

\begin{remark}[On Modifying These Axioms]
Changes to any axiom here propagate throughout the entire theory.
Do not modify Axioms~\ref{ax:algebra}--\ref{ax:fieldeq} without explicit author direction.
See \texttt{THEORY/README.md} for the full list of disclaimers and out-of-scope claims.
\end{remark}

\end{document}
